\documentclass[11pt]{article}
\usepackage[margin=1in]{geometry}
\usepackage{amsmath}
\usepackage{graphicx}
\usepackage{tikz-cd}
\usepackage{multicol}
\usepackage{array}
\usepackage{float}
\usepackage{booktabs}
\usepackage{hyperref}

\setlength{\parindent}{0pt}
\setlength{\parskip}{1\baselineskip}


\begin{document}
\title{Natural Computing - Final Report\\ Emergence of Predator–Prey Strategies via Neuroevolution}
\author{Beatriz Ferrão, Claire Murphy, Yannic Talavera}
% names in alphabetical order ;)
\date{}
\maketitle

% Aim for maximum 10 pages

\section{Problem Statement}

Can predator-prey population dynamics emerge in a neuroevolutionary agent-based simulation
without hard-coded ecological strategy, and under what parameter conditions are these dynamics
most reliably produced?

We reviewed the literature on predator-prey dynamics and identified three consistent
behaviours as success criteria:

\begin{itemize}
    \item \textbf{Population oscillation.} Predator and prey population counts oscillate
          anti-phase, as predicted by the classical Lotka-Volterra
          model~\cite{lotka1925,volterra1926}.
    \item \textbf{Phase-lag relationship.} Prey population peaks consistently precede predator
          peaks—the empirical hallmark of consumer-resource cycles.
    \item \textbf{Red Queen dynamics.} Both species persist across generations and realised
          fitness continues to fluctuate rather than converging to a fixed-point equilibrium,
          indicating ongoing coevolutionary pressure rather than
          stasis~\cite{valen1973,dieckmann1995}.
\end{itemize}

\section{Related Work}



\section{Approach}

\subsection{Simulation Platform}

We use \textit{biosim4}~\cite{biosim4}, an open-source agent-based neuroevolution simulator.
Each simulation runs on a $128 \times 128$ discrete grid populated by 500 agents.
Time is discretized into \emph{generations}, each comprising 500 simulation steps; all
experiments run for 400 generations.
Within each generation, agents sense their local environment through a fixed sensor set, pass
signals through an evolved neural network, and execute the resulting action each simulation
step (subject to per-species action-period gates, described below).
Agents that survive to the end of the generation are eligible for reproduction; offspring
inherit a mutated copy of the parent genome.

\subsubsection{Genome and Neural Architecture}

Each agent's neural network is encoded by a variable-length genome (initial length 14 genes;
maximum 300).
Genes specify directed connections between sensory inputs, hidden neurons, and motor outputs;
network topology and connection weights are therefore evolved jointly.
The number of hidden neurons is bounded by \texttt{maxNumberNeurons} (default 6).
Genetic operators include point mutation (\texttt{pointMutationRate} $= 0.001$) and gene
insertion/deletion (\texttt{geneInsertionDeletionRate} $= 0.0001$).
No crossover is applied; reproduction is asexual.

\subsection{Two-Species Architecture}

Setting \texttt{predatorPreyEnabled = true} splits the population into two co-evolving
species.
A fixed fraction \texttt{predatorFraction} of the 500 agents are initialised as predators each
generation (default 0.30); the remainder are prey.

\textbf{Prey fitness} is the fraction of the generation survived:
$f_{\mathrm{prey}} = \text{age} / \texttt{stepsPerGeneration}$.
Prey are therefore selected jointly by environmental-challenge survival and predator avoidance.

\textbf{Predator fitness} is determined by hunting success:
$f_{\mathrm{pred}} = \min\!\bigl(1,\; \text{captures} / \texttt{predatorCaptureNorm}\bigr)$,
with \texttt{predatorCaptureNorm} $= 5$.
A predator must also make at least \texttt{predatorMinCapturesToReproduce} captures to be
reproduction-eligible (default 3), enforcing a minimum hunting-competence threshold.

Because a capture simultaneously increments predator score and terminates a prey agent, the
two fitness series are near zero-sum by construction (empirically, Pearson $r \approx -0.97$
between generation-mean fitnesses across runs).
This is the reason why the classical criterion of ``both fitnesses trending upward'' is
structurally unreachable in this simulator and would be a misleading success criterion.

\textbf{Starvation feedback.}
A predator that fails to make a kill within \texttt{predatorStarvationSteps} consecutive
simulation steps dies mid-generation.
This couples predator survival to prey abundance: when prey become scarce (because they have
evolved better evasion), more predators starve, relaxing hunting pressure and allowing the prey
population to recover—the negative feedback term that sustains Lotka-Volterra
oscillations~\cite{lotka1925,volterra1926}.

\textbf{Action periods.}
Each species evaluates its neural network and executes an action every $N$th simulation step,
where $N$ is the species-specific action period (\texttt{predatorActionPeriod},
\texttt{preyActionPeriod}).
A period of~1 means the species acts at every step; a period of~2 means it acts every other
step, and so on.
Age and starvation timers tick every step regardless of action period, so larger periods reduce
both the number of decisions taken per generation and the window available for hunting or
evasion.

\subsection{Key Experimental Parameters}

Table~\ref{tab:params} lists the eight parameters varied across our experiments.
Each was chosen to probe a distinct hypothesis about the conditions necessary for
predator-prey coevolution.

\begin{table}[H]
\centering
\caption{Experimental parameters and their roles.}
\label{tab:params}
\begin{tabular}{llll}
\toprule
Parameter & Default & Range swept & Role \\
\midrule
\texttt{challenge}                       & 20    & \{20, 0, 5, 11\}   & Environmental geometry \\
\texttt{predatorFraction}                & 0.30  & 0.10--0.50         & Population balance \\
\texttt{predatorActionPeriod}            & 1     & 1--3               & Predator action speed \\
\texttt{preyActionPeriod}                & 1     & 1--3               & Prey action speed \\
\texttt{predatorStarvationSteps}         & 200   & 0--400             & Ecological feedback strength \\
\texttt{pointMutationRate}               & 0.001 & $10^{-4}$--$10^{-3}$ & Evolutionary exploration rate \\
\texttt{maxNumberNeurons}                & 6     & 2--8               & Neural network capacity \\
\texttt{predatorMinCapturesToReproduce}  & 3     & 0--5               & Reproductive selection pressure \\
\bottomrule
\end{tabular}
\end{table}

\subsection{Statistical Validation}

We use three complementary tests, each targeting a distinct empirical signature of
predator-prey coevolution.
All tests operate on per-generation epoch logs (predator and prey survivor counts and
mean fitness values).
The first 50 generations are excluded as a burn-in to discard bootstrap noise before
meaningful evolutionary dynamics emerge.

\textbf{OSC1 — Population Oscillation.}
The Lotka-Volterra model predicts that predator and prey population counts oscillate
anti-phase~\cite{lotka1925,volterra1926}.
We test this by computing the autocorrelation of the linearly detrended prey survivor series;
linear detrending removes any long-run population drift so that only the oscillatory component
is assessed.
OSC1 passes if the first negative autocorrelation peak satisfies $r < -0.15$ at lag
$\geq 3$ generations.

\textbf{OSC2 — Phase Lag.}
Consumer-resource cycles predict that prey population peaks systematically \emph{precede}
predator peaks.
We compute the cross-correlation of the detrended predator and prey survivor series and
identify the peak correlation at positive lags (prey leads predator).
Statistical significance is assessed by a permutation test (1000 block permutations of the
predator series) at $\alpha = 0.10$.
OSC2 passes if the peak lag is $\geq 2$ generations and the permutation $p$-value is below the
threshold.
This test is the headline diagnostic for oscillatory coevolution: it is robust to run-by-run
variation in cycle amplitude and does not require the oscillations to be large enough to
trigger OSC1.

\textbf{RED-QUEEN — Non-Stasis.}
Following the three-outcome framing of Dieckmann, Marrow, and Law~\cite{dieckmann1995},
any predator-prey system converges to one of three regimes: extinction of one species,
convergence to a fixed-point ESS (evolutionary stasis), or persistent Red Queen cycling.
Because realised fitness is zero-sum, requiring a directional upward trend in both fitnesses
simultaneously—as in the original Van Valen formulation~\cite{valen1973}—is
structurally impossible.
Instead, we test for the third outcome by ruling out the first two.
The test passes if and only if:
\begin{enumerate}
  \item \textbf{Persistence:} both species are present in the final quartile of generations
        (mean survivors $\geq 20\%$ of the per-species early-run baseline); and
  \item \textbf{Non-stasis:} the standard deviation of the linearly detrended realised fitness
        series over the final quartile exceeds 0.02 for at least one species.
\end{enumerate}
Calibration: all no-predation control runs produce a late-window fitness SD of exactly 0.000
(fitness is pinned at a ceiling when predation is absent), so they fail as required.
Treatment runs with active predation produce late-window SDs of 0.02--0.33, confirming the
test is reachable under realistic coevolutionary conditions.

\subsection{Experimental Design}

All experiments share the baseline configuration described above.
Each treatment group varies one or two parameters at a time.
Two no-predation control groups (\texttt{predatorPreyEnabled = false}) are included to verify
that all three statistical tests fail when predation is absent, confirming that any passing
result in a treatment group is caused by predation rather than an artefact of the simulation.

\textbf{Control (CTRL1, CTRL2).}
CTRL1 sweeps the same four environments as B1 with predation disabled.
CTRL2 uses the open arena with predation disabled and starvation turned off
(\texttt{predatorStarvationSteps = 0}).
All three tests must fail on every control run for the experimental results to be
causally interpretable.

\textbf{B: Environment experiments.}
B1 tests whether predator-prey dynamics emerge across four arena geometries—open arena
(\texttt{challenge = 20}), circle (\texttt{challenge = 0}), corner
(\texttt{challenge = 5}), and wall (\texttt{challenge = 11})—to isolate the effect of
prey environmental pressure on coevolution.
B2 crosses mutation rate with two environments to examine whether evolutionary exploration
rate interacts with geometric difficulty.
B3 sweeps predator fraction (0.10--0.50) to identify the population balance that sustains
stable coevolution.

\textbf{Y: Speed experiments.}
Y1 sweeps five predator/prey action-period combinations in the open arena—balanced (1,1),
predator fast (1,2), predator very fast (1,3), predator slow/prey fast (2,1), and
predator slow (2,3)—to map how speed asymmetry affects coevolutionary dynamics.
Y2 extends four representative speed configurations (predator very fast, predator fast,
balanced, predator slow) across two environments; Y3 crosses them with mutation rate to
assess whether speed effects interact with evolutionary exploration.

\textbf{C: Selection pressure experiments.}
C1 sweeps predator starvation pressure (0--400 steps) to identify the feedback strength that
produces Lotka-Volterra-style population regulation.
C2 sweeps the reproduction threshold (0--5 kills required to breed) to test whether stronger
hunting-competence selection drives an arms race.
C3 varies maximum hidden neurons (2--8) to assess whether network expressiveness limits the
emergence of hunting and evasion strategies.
C4 sweeps predator perception radius (3--16 world units) to test whether the range at which
predators can sense prey—a pure predator-side advantage—shifts the coevolutionary balance.


\section{Results}
\subsection{Starvation Pressure as Ecological Feedback} % Claire (C1)
\subsection{Reproductive Gating}                        % Claire (C2)
\subsection{Neural Network Capacity}                    % Claire (C3)
\subsection{Predator Perception Radius}                 % Claire (C4)
\subsection{Effect of Environment Geometry}             % Bia (B1)

\textbf{Question:} Does the spatial structure of the arena affect whether predator-prey
coevolution emerges?

We ran four arena geometries — open arena, circle, corner, and wall — each with 3 independent
random seeds (12 runs total, 400 generations each).
In each geometry, prey must solve a different survival challenge: staying inside a circular
boundary, surviving in a corner region, avoiding a wall, or simply outlasting predators in an
unconstrained space.
Predators always have the same goal regardless of geometry: capture enough prey to reproduce.

\begin{table}[H]
\centering
\caption{B1 results: how many of 3 independent seeds showed each coevolution signature per
environment. \emph{Prey leads predator} = prey population peaks first, then predator peaks
follow (classic boom-bust cycle). \emph{Red Queen} = both species survive and fitness keeps
fluctuating rather than freezing.}
\label{tab:b1}
\begin{tabular}{lcccc}
\toprule
Environment & \shortstack{Population\\oscillation\\(OSC1)} & \shortstack{Prey leads\\predator\\(OSC2)} & \shortstack{Red Queen\\non-stasis\\(RQ)} & \shortstack{Mean prey\\alive per gen} \\
\midrule
Circle  & 1/3 & \textbf{3/3} & \textbf{3/3} & 25.9 \\
Corner  & 0/3 & \textbf{3/3} & \textbf{3/3} & 56.1 \\
Wall    & 0/3 & \textbf{3/3} & 2/3          & 78.2 \\
Open    & 1/3 & 2/3          & 2/3          & 88.1 \\
\bottomrule
\end{tabular}
\end{table}

\begin{figure}[H]
    \centering
    % Replace path below with figures/B1_overview.png after copying the file
    \includegraphics[width=\textwidth]{Output/Sweep/run_20260527_084108/overview/Bia__B1_environment.png}
    \caption{B1 sweep overview across four environments, 3 seeds each.
    \textbf{(A)} Pass rates per environment and test: green = all seeds pass, red = none pass.
    \textbf{(B)} Cross-correlation between predator and prey population counts over time.
    A trough at lag\,=\,0 means the two populations are anti-correlated at any given moment
    (when one is high the other is low); a positive peak at a positive lag means prey peaked
    first, predator followed — the classic predator-prey boom-bust signature.
    The wall environment (green) shows the strongest signal.
    \textbf{(C)} The lag at which predator and prey counts are most correlated, per seed.
    Positive values confirm prey peaks before predator peaks.
    \textbf{(D)} How much prey fitness is still fluctuating in the final quarter of the run.
    Bars above the dashed line indicate the system is still evolving rather than frozen at
    an equilibrium (Red Queen dynamics).}
    \label{fig:b1_overview}
\end{figure}

\textbf{Spatial structure is the key enabler of coevolution.}
The circle and corner environments produced the clearest coevolution signal: in all 3 seeds,
prey population peaks consistently preceded predator peaks (OSC2), and fitness kept
fluctuating throughout the run rather than settling to an equilibrium (Red Queen).
This is likely because constrained arenas force prey into predictable locations, increasing
predator encounter rates and creating a tighter feedback loop: more prey $\to$ predators
thrive $\to$ prey decline $\to$ predators decline $\to$ prey recover.

\textbf{The open arena is less reliable.}
Without spatial constraints, prey can spread across the full $128 \times 128$ grid, reducing
encounter rates and weakening the ecological coupling.
OSC2 passes in only 2 of 3 seeds, and in one seed (seed 3) the predator population actually
peaks \emph{before} prey — the cycle is inverted — suggesting the coevolutionary dynamic
broke down in that run.

\textbf{The wall environment produces extreme population imbalances.}
In the wall geometry, predators struggled to hunt efficiently in a narrow spatial band,
leading to very few predator survivors per generation ($\sim 1.6$--$13.1$) relative to prey
($\sim 8$--$120$ depending on the seed).
The phase relationship (OSC2) still passes consistently, but the Red Queen test fails in one
seed because the predator population becomes so small that its fitness signal essentially
freezes — the system is technically coevolving but predators are barely present.

\textbf{Large-amplitude population oscillations (OSC1) are rare.}
The classic Lotka-Volterra prediction of clearly periodic boom-bust cycles (measurable via
autocorrelation) passes in at most 1 of 3 seeds in any environment.
This does not mean coevolution is absent — the phase relationship (OSC2) and fitness
fluctuation (Red Queen) both pass reliably — but that oscillations in this simulator are
damped rather than large and regular.
400 generations may simply not be long enough for the oscillations to build the amplitude
needed to register in an autocorrelation test.

\subsection{Mutation Rate and Environment}              % Bia (B2)
\subsection{Predator-Prey Population Balance}           % Bia (B3)
\subsection{Speed Asymmetry: Full Grid}                 % Yannic (Y1)
\subsection{Speed Asymmetry Across Environments}        % Yannic (Y2)
\subsection{Speed Asymmetry and Mutation Rate}          % Yannic (Y3)

\section{Discussion}

\section{Conclusion}

\section{Author Contributions}

\section{Code}


\newpage
\begin{thebibliography}{9}

\bibitem{biosim4}
Cox, D.
\textit{biosim4: An Agent-Based Neuroevolution Simulator}.
GitHub repository. \url{https://github.com/davidrmh/biosim4}, 2021.

\bibitem{lotka1925}
Lotka, A.~J.
\textit{Elements of Physical Biology}.
Williams \& Wilkins, Baltimore, 1925.

\bibitem{volterra1926}
Volterra, V.
Fluctuations in the abundance of a species considered mathematically.
\textit{Nature}, 118:558--560, 1926.

\bibitem{valen1973}
Van Valen, L.
A new evolutionary law.
\textit{Evolutionary Theory}, 1:1--30, 1973.

\bibitem{dieckmann1995}
Dieckmann, U., Marrow, P., and Law, R.
Evolutionary cycling in predator-prey interactions: population dynamics and the Red Queen.
\textit{Journal of Theoretical Biology}, 176(1):91--102, 1995.

\bibitem{cliff1995}
Cliff, D. and Miller, G.~F.
Tracking the red queen: Measurements of adaptive progress in co-evolutionary simulations.
In \textit{Advances in Artificial Life (ECAL 1995)}, Lecture Notes in Computer Science,
vol.\ 929, pp.\ 200--218. Springer, 1995.

\bibitem{chen2020predator}
Chen, Y. and Gao, Y.
\textit{Evolving Predator--Prey Agents Using Neuroevolution in Simulated and Real-World
Environments}.
arXiv preprint arXiv:2010.15792, 2020.

\bibitem{jim2000talking}
Jim, K.-C. and Giles, C.~L.
Talking helps: Evolving communicating agents for the predator--prey pursuit problem.
\textit{Artificial Life}, 6(3):237--254, 2000.

\bibitem{paper3}
Yong, C.~H. and Miikkulainen, R.
Coevolution of role-based cooperation in multi-agent systems.
\textit{IEEE Transactions on Autonomous Mental Development}, 1(3), 2009.

\end{thebibliography}

\end{document}
